\chapter{مروری بر کار‌های گذشته}
\label{chp:related_works}
\section{مقدمه}
در طول زمان، روش‌ها، مدل‌ها و الگوریتم‌های متفاوتی جهت حل مسئله ابرتفکیک ارائه شده است. در اینچا سعی می‌شود با انتخاب سه مورد مناسب از تاریخچه تطور موضوع ابرتفکیک تک‌تصویر نگاه خوبی نسبت به روند پیشرفت پژوهش‌ها در این زمینه شکل گیرد. در ابتدا یک مدل نسبتا کلاسیک که با یک روش خلاقانه بر اساس بازنمایی تنک سعی می‌کند مسئله را حل کند بررسی میکنیم. سپس به شرح یک مدل ترکیبی می‌پردازیم که روش بدیع آموزش به ازای هر تصویر را با شبکه‌های عمیق کانولوشنی ترکیب کرده است. در نهایت به سراغ یکی از مدرن‌ترین مقالات در زمینه مسئله ابرتفکیک می‌رویم که از روش مولد-تشخیص‌دهنده برای تولید تصاویر کیفیت‌بالا استفاده می‌کند. هر یک از این پژوهش‌ها نقطه نظر خاصی به موضوع تحقیق داشته‌اند که بررسی آنها را ارزشمند می‌کند.

\section{ابرتفکیک با کدگذاری تنک}
\cite{yang2010scsr} نخستین چارچوب تنکی را برای ابرتفکیک پیشنهاد دادند
که در آن هر پچ تصویر با وضوح پایین و بالایی به‌صورت ترکیبی خطی از
تعداد کمی اتم‌های یک دیکشنری توصیف می‌شود.
آنها نشان دادند که ضرایب تنکی به‌دست‌آمده از فضای وضوح پایین می‌توانند
مستقیم برای بازسازی فضای وضوح بالا به‌کار روند. این ایده پایه‌ای
برای بسیاری از پژوهش‌های پس از آن شد؛ هرچند حل بهینه‌سازی تنک برای
تمام پچ‌ها هزینهٔ محاسباتی بالایی دارد.

\subsection{یادگیری دیکشنری تنک}
\label{subsec:dict_learning}

در روش پیشنهادی \cite{yang2010scsr}، پیشین تنکی بر این مبناست که
هر جفت پچ با وضوح بالا و پایین، همان ترکیب خطی تنک را نسبت به
دو دیکشنری متناظر $\mathbf{D}_{h}$ (بالا) و $\mathbf{D}_{l}$ (پایین)
به‌دست می‌آورد. اگر دو دیکشنری به‌صورت مستقیم از پچ‌پچ‌های نمونه‌برداری
شده استخراج شوند، مطابقت بین پچ‌های متناظر حفظ می‌شود؛ اما این کار
دیکشنری‌های بسیار بزرگ و محاسبهٔ گران‌قیمت ایجاد می‌کند. در این
بخش دو روش برای یادگیری دیکشنری‌های فشرده که همچنان این
هم‌نگاشت را تضمین می‌کند مرور می‌شود.

برای یادگیری دیکشنری تک‌گانه، در ابتدا مسئلهٔ \emph{کدگذاری تنک} به‌صورت زیر فرموله می‌شود
\begin{equation}
	\label{eq:sparse_coding}
	\min_{\mathbf{D},\mathbf{A}}\;
	\underbrace{\|\mathbf{X}-\mathbf{D}\mathbf{A}\|_{F}^{2}}_{\text{بازسازی}}
	\;+\;
	\underbrace{\lambda\|\mathbf{A}\|_{1}}_{\text{تنک‌گری}},
\end{equation}
که $\mathbf{X}=[\mathbf{x}_{1},\dots,\mathbf{x}_{N}]$ مجموعهٔ
پچ‌های آموزشی، $\mathbf{D}$ دیکشنری بیش‌پوشانی و $\mathbf{A}$
ماتریس ضرایب تنک است. برای رفع ابهام مقیاس، ستون‌های $\mathbf{D}$
با $\|\mathbf{d}_{i}\|_{2}=1$ نرمال می‌شوند.

به‌دلیل یادگیری همزمان $\,\mathbf{D}\,$ و $\,\mathbf{A}\,$،
بهینه‌سازی به‌صورت متناوب انجام می‌شود:
\bigskip   % کمی فضای عمودی برای جداسازی

\noindent\textbf{1)} \textbf{مقداردهی اولیه:}\\
\noindent
$\mathbf{D}^{(0)}$ به‌صورت یک ماتریس گوسی تصادفی با ستون‌های نرمال شده تولید می‌شود.
\vspace{0.6\baselineskip}

\noindent\textbf{2)} \textbf{به‌روزرسانی ضرایب:}\\
\noindent
با ثابت نگه داشتن $\mathbf{D}$، مسئله \eqref{eq:sparse_coding}
به‌صورت برنامهٔ خطی \LTRfootnote{Linear Programming} حل می‌شود
\vspace{0.6\baselineskip}

\noindent\textbf{3)} \textbf{به‌روزرسانی دیکشنری:}\\
\noindent
با ثابت نگه داشتن $\mathbf{A}$، یک مسألهٔ QCQP به‌دست می‌آید که با ابزارهای استاندارد
(مثلاً \texttt{CVX}) قابل حل است.
\vspace{0.6\baselineskip}

\noindent\textbf{4)} \textbf{تکرار:}\\
\noindent
گام‌های ۲ و ۳ تا همگرایی $\mathbf{D}$ و $\mathbf{A}$ تکرار می‌شود.

\bigskip

\begin{equation}\label{eq:sparse_coding}
	\min_{\mathbf{A}}\;\|\mathbf{X}-\mathbf{D}\mathbf{A}\|_{F}^{2}
	\quad\text{s.t.}\quad \|\mathbf{a}_{i}\|_{0}\le s .
\end{equation}
در مقاله \cite{yang2010scsr}، برای تصویرهای عمومی از $100\,000$
جفت پچ $(9\times9)$ استفاده شد و دیکشنری نهایی شامل $512$ اتم شد
(شکل ۲).

در گام بعدی هدف این است که دو دیکشنری $\mathbf{D}_{h}\in\mathbb{R}^{m_{h}\times K}$
و $\mathbf{D}_{l}\in\mathbb{R}^{m_{l}\times K}$ طوری آموزش داده شوند
که هر پچ بالایی $\mathbf{x}_{i}^{h}$ و پچ پایین‌دستی متناظر
$\mathbf{x}_{i}^{l}$ یک \emph{کد تنک مشترک}
$\boldsymbol{\alpha}_{i}$ داشته باشند:
\begin{align}
	\min_{\mathbf{D}_{h},\mathbf{D}_{l},\{\boldsymbol{\alpha}_{i}\}} &
	\sum_{i=1}^{N}\Bigl(
	\|\mathbf{x}_{i}^{h}-\mathbf{D}_{h}\boldsymbol{\alpha}_{i}\|_{2}^{2}
	+\|\mathbf{x}_{i}^{l}-\mathbf{D}_{l}\boldsymbol{\alpha}_{i}\|_{2}^{2}
	\Bigr)
	\;+\;
	\lambda\sum_{i=1}^{N}\|\boldsymbol{\alpha}_{i}\|_{1}.
	\label{eq:joint_obj}
\end{align}
در معادله \eqref{eq:joint_obj}، $\lambda$ وزن تعادل بین خطای بازسازی
و تنک‌گری است. با ترکیب دو جملهٔ بازسازی می‌توان معادله را به شکل زیر
بازنویسی کرد
\[
\min_{\mathbf{D},\mathbf{A}}
\bigl\|\mathbf{X}-\mathbf{D}\mathbf{A}\bigr\|_{F}^{2}
\;+\;
\lambda\|\mathbf{A}\|_{1},
\qquad
\mathbf{D}=
\begin{bmatrix}
	\mathbf{D}_{h}\\[2pt]
	\mathbf{D}_{l}
\end{bmatrix},
\;\;
\mathbf{X}=
\begin{bmatrix}
	\mathbf{X}_{h}\\[2pt]
	\mathbf{X}_{l}
\end{bmatrix},
\tag{24}
\label{eq:joint_concat}
\]
که \eqref{eq:joint_concat} دقیقاً همان فرم \eqref{eq:sparse_coding}
است؛ به‌جای یک دیکشنری واحد، یک دیکشنری مسدود \LTRfootnote{block‑diagonal}
ساخته می‌شود. به همین دلیل، همان استراتژی متناوب الگوریتم
تک‌دیکشنری می‌تواند برای $\mathbf{D}_{h}$ و $\mathbf{D}_{l}$
به‌صورت همزمان به‌کار رود. نتایج نشان می‌دهد دیکشنری‌های
یادگرفته‌شده الگوهای پایه‌ای مانند لبه‌های جهت‌دار، خطوط و بافت
های ساده را نشان می‌دهند؛ برخلاف نمونهٔ خام پچ‌ها که انعطاف‌پذیری
کمتری دارند.
	\begin{figure}
	\centering
	\includegraphics[width=1\linewidth]{figures/scsr_overall.png}
	\caption{تصاویر از بالا چپ به سمت پایین راست به ترتیب مربوط به ورودی کم کیفیت، درونیابی دومکعبی، بازافکنی، NE, SE, ScSR می‌باشند.}
	\label{fig:scsr_overall}
	\end{figure}
	
	\subsection{نتایج تجربی}\label{sec:exp_results}
	در این بخش عملکرد الگوریتم پیشنهادی بر روی دو دستهٔ تصویر \textit{عمومی} و \textit{چهره} ارزیابی می‌شود و تأثیر پارامترهای کلیدی اندازهٔ دیکشنری، نویزو قیدهای بازسازی سراسری مورد بررسی قرار می‌گیرد. تمامی آزمایش‌ها بر مبنای بزرگ‌نمایی ورودی با نسبت ۳× برای تصاویر عمومی و ۴× برای تصاویر چهره انجام شده است؛ این مقادیر در مطالعات پیشین \cite{GlasnerSR2009} نیز رایج هستند.
	
در این زیرمجموعه دو آزمایش مجزا برای \textbf{تصاویر عمومی} و \textbf{تصاویر چهره} گزارش می‌شود.
برای تصاویر عمومی داریم:
	\begin{itemize}
		\item \textbf{پچ‌ها:} برای تصاویر عمومی پچ‌های $3\times3$ (رزولوشن پایین) که به $6\times6$ ارتقاء می‌یابند، با هم‌پوشانی $1$ پیکسل استفاده می‌شود؛ معادل این موارد در رزولوشن بالا $9\times9$ با هم‌پوشانی $3$ پیکسل است.
		\item \textbf{دیکشنری‌ها:} دو دیکشنری \(\mathbf{D}_{h}\) و \(\mathbf{D}_{l}\) با اندازهٔ پیش‌فرض $K=1024$ از $10^{5}$ جفت پچ طبیعی استخراج شدند (فقط نواحی بافت‌دار نگهداری شد) \cite{yang2010scsr}.
		\item \textbf{پارامتر منظم‌سازی:} برای تصاویر کم‌نویز مقدار $\lambda=0.1$ به‌کار رفت؛ مقدار بزرگتر $\lambda$ برای ورودی‌های پرنویز انتخاب می‌شود.
	\end{itemize}
	تمامی معیارهای عددی بر \textit{کانال روشنایی} (Y) محاسبه شده‌اند؛ همان‌طور که اشاره شده، RMSE تنها یک معیار کمّی است و با ارزیابی بصری کاملاً هم‌راستا نیست.
	
	\paragraph
در \ref{tbl:gaussian_noise_ablation} نتایج ما در کنار \textit{Bicubic} و \textit{Neighborhood Embedding (NE)} \cite{ChangNeighbor2004} نشان داده شده‌اند.
	
	\begin{table}[!ht]
		\centering
		\caption{تأثیر سطوح مختلف نویز گاوسی بر روی روش‌‌های مختلف}
		\label{tbl:gaussian_noise_ablation}
		% فاصلهٔ خطوط ردیف‌ها
		\renewcommand{\arraystretch}{1.4}
		%-----------------------------------------------------------------
		\begin{tabular}{|c|c|c|c|c|}
			\hline
			\textbf{روش / سطح نویز} & \textbf{0} & \textbf{4} & \textbf{6} & \textbf{8} \\ \hline
			دومکعبی                 & \lr{9.873}  & \lr{10.423} & \lr{11.037} & \lr{11.772} \\ \hline
			همسایگی                  & \lr{9.534}  & \lr{10.734} & \lr{11.856} & \lr{13.064} \\ \hline
			نمایش تنک                & \lr{8.359}  & \lr{9.240}  & \lr{10.184} & \lr{11.448} \\ \hline
		\end{tabular}
		%-----------------------------------------------------------------
	\end{table}
	
	\begin{enumerate}
		\item \textbf{حفظ بافت:} حاشیهٔ برگ‌ها و لکه‌های صورت در \textit{NE} به‌طور واضحی خاکستری هستند؛ الگوریتم ما جزئیات ریزتر (مانند موی شاخهٔ برگ) را بازمی‌گرداند.
		\item \textbf{دقت عددی:} RMSE برای روش پیشنهادی $0.87$ (تصویر «Flower») بوده که از \textit{Bicubic} ($1.02$) و \textit{NE} ($0.94$) بهتر است. نکتهٔ مهم این است که هرچند \textit{NE} از نظر بصری بهتر از \textit{Bicubic} به‌نظر می‌رسد، RMSE آن بالاتر است؛ همان‌طور که پیش‌تر اشاره شد، RMSE معیار کاملی برای کیفیت بصری نیست.
	\end{enumerate}
	
	\subsection{تصاویر چهره}
پروتکل آزمایش‌ها در مجموعهٔ \textit{FRGC} به‌صورت زیر بود:
	
	\begin{itemize}
		\item تصاویر HR با اندازه $100\times100$ پیکسل قبلاً با الگوریتم خودکار هم‌ترازی چشم‌ها تنظیم شدند.
		\item برای ساخت جفت پچ‌ها، $10^{5}$ جفت $5\times5$ از 540 تصویر آموزشی استخراج شد و دو دیکشنری $K=1024$ آموخته شد.
		\item تصاویر LR به $25\times25$ پیکسل کاهش یافته (نویز گاوسی $\sigma=5$) و پس از هم‌ترازی دستی، به‌عنوان ورودی الگوریتم استفاده شد.
	\end{itemize}
	
	\paragraph{رویکرد دو‑مرحله‌ای.}
	روش پیشنهادی شامل دو گام است:
	\begin{enumerate}
		\item \textbf{مدل سراسری NMF:} بازسازی اولیه با منظم‌سازی $\lambda=0.3$.
		\item \textbf{مدل موضعی تنکی:} بهبود حاشیه‌ها و بافت‌های دقیق بر پایهٔ \eqref{eq:joint_obj}.
	\end{enumerate}
	
	\begin{figure}
		\centering
		\includegraphics[width=0.7\linewidth]{figures/ScSR_Faces.png}
		\caption{ستون a تصویر کم کیفیت، ستون b حاصل درونیابی دومکعبی، ستون c حاصل بازافکنی و d حاصل NMF سپس فیلتر دوطرقه و e حاصل NMF و سپس دیکشنری تنک  و f تصویر اصلی است. }
		\label{fig:ScSR_Faces}
	\end{figure}
	
	همانطور که مشاهده میکنید در \ref{fig:ScSR_Faces} مقدار RMSE برای ستونهای b, c, d, e به ترتیب مقادیر ۸/۰۲۴، ۷/۴۲۴، ۱۰/۷۳۸ و ۶/۸۹۱ می‌باشد.
	
	\subsection{تاثیر اندازه دیکشنری}
	به منظور بررسی تعادل میان توان بازنمایی و هزینهٔ محاسباتی، چهار دیکشنری با $K\in\{256,512,1024,2048\}$ آموزش داده شد. همچنین مجموعهٔ کامل $10^{5}$ پچ به‌عنوان دیکشنری \textit{غیر‌فشرده} (مبنای \cite{yang2010scsr}) مورد استفاده قرار گرفت.
	
	\begin{table}[htbp]
		\centering
		\caption{RMSE بازسازی تصویر \lr{Lena} با اندازه‌های مختلف دیکشنری (پیکسل).}
		\label{tab:dict_rmse}
		\begin{tabular}{c|c}
			\toprule
			$K$ & RMSE \\
			\midrule
			256   & 1.47 \\
			512   & 1.29 \\
			1024  & \textbf{1.12} \\
			2048  & 1.09 \\
			$10^{5}$ (دیکشنری غیرفشرده) & 1.07 \\
			Bicubic & 1.53 \\
			\bottomrule
		\end{tabular}
	\end{table}
	
	همان‌طور که در \ref{fig:Dic_Performance} مشاهده می‌شود، افزایش $K$ به‌تدریج علائم آرشیو (الگوهای نادر) را بهتر بازپیکربندی می‌کند؛ با این حال، زمان محاسبه به‌صورت تقریباً خطی با $K$ رشد می‌کند (جدول \ref{tab:dict_time}).
	
	\begin{table}[htbp]
		\centering
		\caption{زمان اجرا (ثانیه) برای تصویر \lr{Girl} بر روی لپ‌تاپ Core Duo @ 1.83 GHz.}
		\label{tab:dict_time}
		\begin{tabular}{c|c}
			\toprule
			$K$ & زمان (s) \\
			\midrule
			256   & 3.1 \\
			512   & 5.4 \\
			1024  & 9.8 \\
			2048  & 18.7 \\
			$10^{5}$ & 86.2 \\
			\bottomrule
		\end{tabular}
	\end{table}
	
	\begin{figure}
		\centering
		\includegraphics[width=0.5\linewidth]{figures/Dic_Performance.png}
		\caption{رابطه نسبتا خطی اندازه دیکشنری و سرعت اجرای الگوریتم}
		\label{fig:Dic_Performance}
	\end{figure}
	
	\textbf{نتیجه‌گیری:} برای اکثر کاربردها $K=1024$ نقطهٔ بهینه‌ای بین کیفیت (RMSE ≈ 1.1) و سرعت (≈ 10 s برای تصویر $256\times256$) ارائه می‌دهد.
	
	\subsection{پایداری نسبت به نویز}\label{subsec:noise}
	اکثراً الگوریتم‌های SR فرض می‌کنند تصویر ورودی تمیز است؛ اما در عمل نویز گاوسی با انحراف استاندارد $\sigma\in[4,10]$ رایج است. در \eqref{eq:joint_obj} پارامتر منظم‌سازی $\lambda$ به‌صورت $\lambda=0.1\,\sigma$ تنظیم شد.
	
	\paragraph{نتایج تصویری.}
	
	روش \textit{NE} می‌تواند لبه‌ها را حفظ کند، اما با افزایش $\sigma$ نشانه‌های نویز را نیز تقویت می‌نماید؛ در مقابل، الگوریتم ما هر دو عمل \textit{دینویزینگ} و \textit{Super‑Resolution} را به‌صورت همزمان انجام می‌دهد و نتایج واضح‌تری ارائه می‌کند.
	
	\begin{table}[htbp]
		\centering
		\caption{RMSE برای سطوح مختلف نویز ($\sigma$) بر روی تصویر \lr{Liberty Statue}.}
		\label{tab:noise_rmse}
		\begin{tabular}{c|ccc}
			\toprule
			$\sigma$ & Bicubic & NE & پیشنهادی \\
			\midrule
			4  & 1.46 & 1.38 & \textbf{1.21} \\
			7  & 1.78 & 1.71 & \textbf{1.42} \\
			10 & 2.03 & 1.97 & \textbf{1.68} \\
			\bottomrule
		\end{tabular}
	\end{table}
	
	\subsection{اثر قیدهای بازسازی سراسری}\label{subsec:global}
	قید سراسری\LTRfootnote{Global Reconstruction Constraint} پس از مرحلهٔ موضعی اعمال می‌شود تا تصویر بازسازی‌شده با دادهٔ LR سازگار باشد. آزمایش‌های زیر نشان می‌دهند که این قید تنها به‌صورت کمی RMSE را کاهش می‌دهد؛ اما در برخی نقاط مرزی، به‌ویژه در مناطق تک‌رنگ، آرتیفکت‌های ناخواسته را از بین می‌برد.
	
	\begin{table}[htbp]
		\centering
		\caption{مقایسه RMSE بین مدل موضعی تنها و ترکیب موضعی + سراسری (دیکشنری $K=1024$).}
		\label{tab:global_rmse}
		\begin{tabular}{c|cc}
			\toprule
			تصویر & فقط موضعی & موضعی + سراسری \\
			\midrule
			Lena   & 1.12 & \textbf{1.08} \\
			Girl   & 1.24 & \textbf{1.19} \\
			Statue & 1.31 & \textbf{1.26} \\
			\bottomrule
		\end{tabular}
	\end{table}
	
	\subsection{خلاصهٔ نتایج}
	\begin{itemize}
		\item استفاده از \textit{دیکشنری‌های جفت‌شده} (بالا/پایین) و ترکیب وزن‌های \textit{تنک} امکان بازسازی جزئیات ریز در هر دو دامنهٔ عمومی و چهره را فراهم می‌کند.
		\item اندازهٔ $K=1024$ بهترین توازن بین کیفیت (RMSE ≈ 1.1) و سرعت (≈ 10 s برای تصویر $256\times256$) را ارائه می‌دهد.
		\item پارامتر منظم‌سازی $\lambda$ به‌عنوان تابعی از میزان نویز ($\lambda\propto\sigma$) به‌صورت خودکار تنظیم می‌شود و موجب ترکیب همزمان \textit{denoising} و \textit{SR} می‌گردد.
		\item قید سراسری، هرچند اثر کمی بر RMSE دارد، اما به‌صورت قابل ملاحظه‌ای اثرات آرتیفکت را در نواحی یکنواخت کاهش می‌دهد.
	\end{itemize}

\section{ابرتفکیک بدون نمونهٔ آموزشی با یادگیری عمیق داخلی}
\label{sec:zssr}

\subsection{مقدمه}
\label{subsec:zssr_intro}

یادگیری عمیق در سال‌های اخیر جهشی چشمگیر در عملکرد ابرتفکیک تصویر
ایجاد کرده است \cite{lim2017enhanceddeepresidualnetworks, wang2018esrganenhancedsuperresolutiongenerative}.
با این حال، روش‌های مبتنی بر یادگیری نظارت‌شده با محدودیتی اساسی
مواجه‌اند: آن‌ها فرض می‌کنند فرایند کاهش وضوح از تصاویر با کیفیت بالا
به تصاویر کم‌کیفیت
از پیش مشخص و ثابت است. به‌طور معمول، هستهٔ
کاهش‌نمونه‌برداری یک هستهٔ دومکعبی
\LTRfootnote{Bicubic kernel}
با ضد‌اعوجاج است (دستور پیش‌فرض \texttt{imresize} نرم‌افزار متلب)
و فرض می‌شود هیچ‌گونه آرتیفکت اضافی مانند نویز حسگر، فشرده‌سازی تصویر
یا تابع انتشار نقطه‌ای\LTRfootnote{Point Spread Function (PSF)} غیرایده‌آل
وجود ندارد \cite{shocher2018zssr}.

در عمل، تصاویر کم‌کیفیت واقعی به‌ندرت از این فرضیات پیروی می‌کنند.
تصاویر دانلود‌شده از اینترنت، عکس‌های گرفته‌شده با تلفن همراه، تصاویر
تاریخی قدیمی و داده‌های تصویری زیستی هر یک فرایند تصویربرداری منحصر
به‌فردی دارند. در چنین شرایطی، روش‌های ابرتفکیک پیشرفته مبتنی بر
یادگیری نظارت‌شده عملکردی ضعیف از خود نشان می‌دهند. به‌عنوان مثال،
شبکهٔ عمیق EDSR \cite{lim2017enhanceddeepresidualnetworks} که بر داده‌های ایده‌آل آموزش دیده،
هنگام مواجهه با تصاویری که هستهٔ کاهش‌نمونه‌برداری آن‌ها غیردومکعبی
است یا آرتیفکت‌های حاصل از نویز و فشرده‌سازی دارند، نتایجی تار و
بدون جزئیات تولید می‌کند.

این مسئله انگیزهٔ اصلی شوخر و همکاران \cite{shocher2018zssr} را برای
ارائهٔ روش \lr{ZSSR}\LTRfootnote{Zero-Shot Super-Resolution} فراهم ساخت.
ایدهٔ کلیدی این پژوهش بهره‌گیری از قدرت یادگیری عمیق است، بدون آنکه
به هیچ دادهٔ آموزشی خارجی یا پیش‌آموزش وابسته باشد. به جای آن، از
تکرارپذیری درونی اطلاعات در یک تصویر واحد استفاده می‌شود تا یک شبکهٔ
عصبی کانولوشنی خاص هر تصویر در زمان آزمون آموزش داده شود. تا جایی که
نویسندگان اطلاع دارند، ZSSR نخستین روش ابرتفکیک مبتنی بر CNN بدون
نظارت است.
\begin{figure}
	\centering
	\includegraphics[width=1\linewidth]{figures/ZSSR_first.png}
	\caption{تصویری تاریخی از پایان جنگ جهانی دوم. در ردیف بالا روش EDSR \cite{lim2017enhanceddeepresidualnetworks} و در ردیف پایین روش ZSSR را مشاهده می‌کنید.}
	\label{fig:ZSSR_first}
\end{figure}

\subsection{قدرت پیش‌بینی آمار درونی تصویر}
\label{subsec:internal_stats}

پایهٔ نظری روش ZSSR بر مشاهده‌ای بنیادین استوار است: تصاویر طبیعی
تکرارپذیری درونی قوی در داده‌ها دارند. پچ‌های کوچک تصویر (مثلاً
$5\times5$ یا $7\times7$) بارها و بارها در درون یک تصویر واحد تکرار
می‌شوند، هم در مقیاس یکسان و هم در مقیاس‌های مختلف
\cite{GlasnerSR2009}.

این مشاهده به‌صورت تجربی توسط گلاسنر و همکاران
\cite{GlasnerSR2009} و زونتاک و ایرانی \cite{zontak2011internal}
بر روی صدها تصویر طبیعی تأیید شده و نشان داده شده که تقریباً برای
هر پچ کوچک و در تقریباً هر تصویر طبیعی صادق است. این ویژگی پایه‌ای
بسیاری از روش‌های بدون نظارت برای بهبود تصویر از جمله ابرتفکیک بدون
نظارت \cite{GlasnerSR2009}،
ابرتفکیک کور \cite{MittalCompletelyBlind} و محو‌زدایی کور
بوده است.

شوخر و همکاران به یکی از مهم‌ترین شواهد این قدرت پیش‌بینی درونی اشاره
می‌کنند: در تصویری از ساختمانی با بالکن‌های متعدد، نرده‌های بسیار
ریز بالکن‌های کوچک تنها از درون خود تصویر قابل بازسازی هستند، زیرا
شواهدی از وجود آن‌ها در بالکن‌های بزرگ‌تر (در مکان و مقیاس دیگر درون
همان تصویر) یافت می‌شود. این شواهد در هیچ مجموعه‌دادهٔ خارجی، هر چقدر
هم بزرگ باشد، وجود ندارد. روش‌های ابرتفکیک پیشرفته‌ای مانند
EDSR \cite{lim2017enhanceddeepresidualnetworks} که بر داده‌های خارجی آموزش دیده‌اند، در
بازسازی این اطلاعات خاصِ تصویر ناموفق‌اند.

علاوه بر این، زونتاک و ایرانی \cite{zontak2011internal} به‌صورت
تجربی نشان دادند که آنتروپی درونی پچ‌ها در یک تصویر واحد
به‌مراتب کمتر از آنتروپی خارجی پچ‌ها در مجموعه‌ای عمومی از
تصاویر طبیعی است. این یافته بدان معناست که آمار درونی تصویر
اغلب قدرت پیش‌بینی قوی‌تری نسبت به آمار خارجی فراهم می‌کند.
این برتری به‌ویژه در شرایط عدم قطعیت بیشتر و تخریب‌های تصویری
شدیدتر تقویت می‌شود.

با وجود مزایای روش‌های بدون نظارت مبتنی بر پچ، آن‌ها با محدودیت‌هایی
مواجه‌اند: معمولاً بر شباهت اقلیدسی ساده میان پچ‌های کوچک با
اندازهٔ ثابت (مثلاً $5\times5$) و جستجوی $K$-نزدیک‌ترین همسایه
تکیه می‌کنند. در نتیجه، توانایی تعمیم‌پذیری آن‌ها برای پچ‌هایی که در
تصویر LR وجود ندارند محدود است و نمی‌توانند معیارهای شباهت ضمنی
یادگیری‌شده یا اندازه‌های غیریکنواخت ساختارهای تکرارشونده در تصویر
را مدیریت کنند.

\subsection{شبکهٔ عصبی کانولوشنی خاص تصویر}
\label{subsec:image_specific_cnn}

روش ZSSR قدرت پیش‌بینی و آنتروپی پایین اطلاعات درونی خاص تصویر را با
قابلیت‌های تعمیم‌پذیری یادگیری عمیق ترکیب می‌کند. فرض کنید تصویر
آزمون $I$ داده شده و هیچ نمونهٔ خارجی برای آموزش در دسترس نیست.
در این حالت، یک شبکهٔ CNN خاص تصویر ساخته می‌شود که منحصراً
برای حل مسئلهٔ ابرتفکیک همین تصویر طراحی شده است.

\subsection{استخراج نمونه‌های آموزشی از تصویر}

ایدهٔ کلیدی آن است که شبکه بر روی نمونه‌هایی آموزش ببیند که صرفاً
از خود تصویر آزمون استخراج شده‌اند. این نمونه‌ها به‌صورت زیر تولید
می‌شوند: تصویر LR ورودی $I$ با ضریب مقیاس مطلوب $s$ کاهش‌نمونه‌برداری
می‌شود تا نسخه‌ای با وضوح پایین‌تر، یعنی $I{\downarrow s}$، به دست آید.
سپس شبکه آموزش می‌بیند تا تصویر $I$ را از نسخهٔ کاهش‌یافته‌اش
$I{\downarrow s}$ بازسازی کند. پس از آموزش، همان شبکه بر روی تصویر
ورودی $I$ اعمال می‌شود و خروجی با وضوح بالا، یعنی $I{\uparrow s}$،
تولید می‌گردد. از آنجا که شبکهٔ آموزش‌دیده کاملاً کانولوشنی است،
قادر به اعمال بر تصاویر با اندازه‌های مختلف خواهد بود.

\subsection{افزایش داده}

از آنجا که مجموعهٔ آموزشی تنها شامل یک نمونه (تصویر آزمون) است، از
روش‌های افزایش داده برای استخراج جفت‌نمونه‌های بیشتر استفاده می‌شود.
تصویر آزمون $I$ به نسخه‌های کوچک‌تر متعددی کاهش‌نمونه‌برداری می‌شود:
$I = I_0, I_1, I_2, \ldots, I_n$. این نسخه‌ها نقش
\emph{پدرهای HR}\LTRfootnote{HR fathers} را ایفا می‌کنند. سپس هر
پدر HR با ضریب مقیاس $s$ کاهش‌نمونه‌برداری می‌شود تا
\emph{فرزندان LR}\LTRfootnote{LR sons} ایجاد شوند. مجموعهٔ حاصل شامل
جفت‌نمونه‌های LR-HR خاص تصویر است.

علاوه بر کاهش مقیاس، هر جفت LR-HR با استفاده از ۴ دوران
($0^\circ, 90^\circ, 180^\circ, 270^\circ$) و بازتاب آینه‌ای در
جهت‌های افقی و عمودی غنی‌سازی می‌شود. این عملیات $8\times$ نمونهٔ
آموزشی خاص تصویر بیشتر اضافه می‌کند.

\subsection{ابرتفکیک تدریجی}

به‌منظور استحکام بیشتر و امکان اعمال ضرایب مقیاس بزرگ حتی بر تصاویر
LR بسیار کوچک، ابرتفکیک به‌صورت تدریجی انجام می‌شود
\cite{GlasnerSR2009}. الگوریتم برای چندین ضریب مقیاس
میانی $(s_1, s_2, \ldots, s_m = s)$ اجرا می‌شود. در هر مقیاس میانی
$s_i$، تصویر SR تولیدشده $\text{HR}_i$ و نسخه‌های
کاهش‌یافته/دوران‌یافتهٔ آن به‌عنوان پدرهای HR جدید به مجموعهٔ
آموزشی (که به‌تدریج در حال رشد است) افزوده می‌شوند. سپس این نمونه‌ها
با ضریب مقیاس بعدی $s_{i+1}$ کاهش‌نمونه‌برداری شده و جفت‌نمونه‌های
آموزشی LR-HR جدید تولید می‌شوند. این فرایند تا رسیدن به افزایش
وضوح کامل مطلوب $s$ ادامه می‌یابد.

\subsection{معماری و بهینه‌سازی شبکه}
\label{subsec:arch_optim}

شبکه‌های CNN نظارت‌شده که بر مجموعه‌ای بزرگ و متنوع از نمونه‌های
خارجی LR-HR آموزش می‌بینند، باید تنوع عظیم تمام روابط ممکن LR-HR
را در وزن‌های یادگیری‌شدهٔ خود رمزگذاری کنند. به همین دلیل، این
شبکه‌ها معمولاً بسیار عمیق و پیچیده هستند. در مقابل، تنوع روابط
LR-HR در یک تصویر واحد به‌مراتب کمتر است و می‌تواند توسط شبکه‌ای
بسیار کوچک‌تر و ساده‌تر رمزگذاری شود.

معماری پیشنهادی ZSSR یک شبکهٔ کاملاً کانولوشنی ساده با مشخصات زیر است:

\begin{itemize}
	\item \textbf{عمق:} ۸ لایهٔ پنهان
	\item \textbf{کانال‌ها:} هر لایه ۶۴ کانال
	\item \textbf{تابع فعال‌سازی:} \lr{ReLU} بر هر لایه
	\item \textbf{ورودی:} تصویر LR درون‌یابی‌شده به اندازهٔ خروجی
	\item \textbf{یادگیری:} تنها باقی‌ماندهٔ\LTRfootnote{Residual} بین
	تصویر LR درون‌یابی‌شده و پدر HR آن یادگیری می‌شود (مشابه
	\cite{zhang2018imagesuperresolutionusingdeep})
\end{itemize}

برای بهینه‌سازی شبکه، از تابع هزینهٔ $L_1$ به‌همراه بهینه‌ساز
\lr{Adam} استفاده شده است. نرخ یادگیری اولیه
$0.001$ در نظر گرفته می‌شود. به‌صورت دوره‌ای، یک برازش خطی بر خطای
بازسازی انجام می‌شود؛ اگر انحراف استاندارد از شیب برازش خطی بیشتر
باشد، نرخ یادگیری بر ۱۰ تقسیم می‌شود. آموزش زمانی متوقف می‌شود که
نرخ یادگیری به $10^{-6}$ برسد. این استراتژی کاهش نرخ یادگیری را
می‌توان به‌صورت زیر خلاصه کرد:
\begin{equation}
	\label{eq:lr_schedule}
	\eta_{t+1} =
	\begin{cases}
		\eta_t / 10, & \text{اگر}\; \sigma_{\text{err}} > c \cdot |\text{slope}| \\
		\eta_t, & \text{در غیر این صورت}
	\end{cases}
\end{equation}
که $\eta_t$ نرخ یادگیری در دورهٔ $t$، $\sigma_{\text{err}}$ انحراف
استاندارد خطای بازسازی و $c$ یک ثابت آستانه‌ای است.
\begin{figure}
	\centering
	\includegraphics[width=1\linewidth]{figures/ZSSR_architecture.png}
	\caption{تفاوت معماری ZSSR با مدل‌های یادگیری بانظارت}
	\label{fig:ZSSR_architecture}
\end{figure}


برای شتاب‌بخشیدن به مرحلهٔ آموزش و استقلال زمان اجرا از اندازهٔ
تصویر آزمون $I$، در هر تکرار یک برش تصادفی با اندازهٔ ثابت (معمولاً
$128 \times 128$) از یک جفت پدر-فرزند انتخابی تصادفی گرفته می‌شود
(مگر آنکه جفت تصویر انتخاب‌شده کوچک‌تر باشد). احتمال نمونه‌برداری
هر جفت LR-HR در هر تکرار آموزشی غیریکنواخت و متناسب با اندازهٔ
پدر HR است: هرچه نسبت اندازهٔ پدر HR به تصویر آزمون $I$ به ۱
نزدیک‌تر باشد، احتمال نمونه‌برداری آن بیشتر است. این طراحی بازتابی
از اعتبار بیشتر نمونه‌های HR غیرمصنوعی نسبت به نمونه‌های سنتز‌شده
است.

در نهایت، از تکنیک \emph{خودتجمیع هندسی}
\LTRfootnote{Geometric self-ensemble} مشابه \cite{lim2017enhanceddeepresidualnetworks}
استفاده می‌شود. ۸ خروجی مختلف برای ۸ دوران و بازتاب تصویر آزمون $I$
تولید و سپس ترکیب می‌شوند. در ZSSR به جای میانگین از \emph{میانه}
استفاده می‌شود. این خروجی همچنین با تکنیک بازافکنی
\LTRfootnote{Back-projection} \cite{GlasnerSR2009}
ترکیب می‌شود: هر ۸ تصویر خروجی چند تکرار بازافکنی را طی کرده
و نهایتاً تصویر میانه نیز با بازافکنی اصلاح می‌گردد.

\paragraph{زمان اجرا.}
هرچند آموزش در زمان آزمون انجام می‌شود، میانگین زمان اجرا برای هر
تصویر ۵۴ ثانیه برای یک افزایش مقیاس واحد (بر GPU مدل \lr{Tesla K-80})
است. این زمان مستقل از اندازهٔ تصویر و ضریب مقیاس نسبی $s$ است (به
دلیل برش‌های تصادفی با اندازهٔ ثابت). با این حال، نتایج بهتر با
افزایش تدریجی وضوح به دست می‌آید. برای مثال، استفاده از ۶ ضریب مقیاس
میانی PSNR را حدوداً $0.2$~dB بهبود می‌بخشد، اما زمان اجرا را به حدود
۵ دقیقه افزایش می‌دهد. مقایسه با EDSR+ \cite{lim2017enhanceddeepresidualnetworks} نشان
می‌دهد که بر همان بستر سخت‌افزاری، EDSR+ حدود ۲۰ ثانیه برای ابرتفکیک
$\times 2$ یک تصویر $200 \times 200$ نیاز دارد، اما زمان اجرای آن
به‌صورت مربعی با اندازهٔ تصویر رشد می‌کند و برای تصویر $800 \times 800$
به ۵ دقیقه می‌رسد. بنابراین، فراتر از این اندازه، ZSSR سریع‌تر از
EDSR+ عمل می‌کند.

\subsection{سازگاری با تصویر آزمون}
\label{subsec:adaptation}

یکی از مزایای بارز ZSSR نسبت به روش‌های نظارت‌شده، قابلیت سازگاری
با تخریب‌ها و تنظیمات خاص هر تصویر آزمون در زمان آزمون است. در
روش‌های نظارت‌شده، شبکه برای یک هستهٔ کاهش‌نمونه‌برداری ثابت (معمولاً
دومکعبی) و شرایط تصویربرداری با کیفیت بالا بهینه شده است. در عمل،
فرایند تصویربرداری از تصویری به تصویر دیگر متفاوت است: دوربین‌ها و
حسگرها متنوع‌اند (لنزها و PSFهای مختلف)، شرایط فردی تصویربرداری
متغیر است (مثلاً لرزش ناخواستهٔ دوربین، شرایط دید ضعیف) و
این تفاوت‌ها به هسته‌های کاهش‌نمونه‌برداری مختلف، ویژگی‌های نویزی
گوناگون و آرتیفکت‌های فشرده‌سازی مختلف منجر می‌شود.

آموزش عملی برای تمام ترکیب‌های ممکن تنظیمات تصویربرداری غیرممکن
است. علاوه بر آن، یک CNN نظارت‌شدهٔ واحد بعید است برای تمام انواع
تخریب‌ها/تنظیمات عملکرد خوبی داشته باشد. دستیابی به عملکرد مطلوب
نیاز به شبکه‌های تخصصی متعدد دارد که هر یک بر نوع خاصی از تخریب‌ها
آموزش دیده باشند (برای روزها یا هفته‌ها). ZSSR با اجازه دادن به
کاربر برای تنظیم پارامترهای زیر در زمان آزمون، این مشکل را حل
می‌کند:

\begin{enumerate}
	\item \textbf{هستهٔ کاهش‌نمونه‌برداری مطلوب:} هستهٔ مورد نظر
	(در صورت عدم ارائه، هستهٔ دومکعبی به‌عنوان پیش‌فرض استفاده می‌شود).
	\item \textbf{ضریب مقیاس ابرتفکیک $s$:} ضریب بزرگ‌نمایی مطلوب.
	\item \textbf{تعداد افزایش‌های تدریجی مقیاس:} تعادل بین سرعت و
	کیفیت (پیش‌فرض ۶).
	\item \textbf{اعمال بازافکنی:} آیا قید سازگاری بین تصاویر LR و HR
	اعمال شود (پیش‌فرض: بله).
	\item \textbf{افزودن نویز به فرزندان LR:} آیا نویز مصنوعی به
	فرزندان LR در هر جفت نمونهٔ آموزشی افزوده شود (پیش‌فرض: خیر).
\end{enumerate}

دو پارامتر آخر (حذف بازافکنی و افزودن نویز) امکان مدیریت ابرتفکیک
تصاویر LR با کیفیت پایین را فراهم می‌سازند. نویسندگان دریافتند که
افزودن مقدار کمی نویز گاوسی (با میانگین صفر و انحراف استاندارد
$\sim 5$ سطح خاکستری) عملکرد را برای طیف وسیعی از تخریب‌ها بهبود
می‌بخشد. علت این پدیده آن است که اطلاعات خاص تصویر تمایل به تکرار
در مقیاس‌های مختلف دارند، در حالی که آرتیفکت‌های نویزی این ویژگی
را ندارند. افزودن نویز مصنوعی به
فرزندان LR (بدون افزودن به پدرهای HR) به شبکه می‌آموزد که
اطلاعات نامرتبط بین‌مقیاسی (نویز) را نادیده بگیرد و تنها وضوح
اطلاعات مرتبط (جزئیات سیگنال) را افزایش دهد.

شایان ذکر است که ارائهٔ هستهٔ کاهش‌نمونه‌برداری تخمین‌زده‌شده به
روش‌های نظارت‌شدهٔ پیشرفته در زمان آزمون سودمند نخواهد بود، زیرا
آن‌ها نیاز به بازآموزش کامل شبکه بر مجموعهٔ جدیدی از جفت‌های
LR-HR دارند که با این هستهٔ خاص (ناپارامتری) تولید شده باشند.

\subsection{آزمایش‌ها و نتایج}
\label{subsec:zssr_experiments}

روش ZSSR عمدتاً برای تصاویر LR واقعی با تنظیمات تصویربرداری ناشناخته
و متغیر طراحی شده است. این تصاویر معمولاً فاقد حقیقت زمینی HR هستند
و ارزیابی بصری روش اصلی سنجش آن‌هاست. با این حال، برای ارزیابی کمّی،
چندین آزمایش کنترل‌شده در تنظیمات مختلف انجام شده است.

\subsubsection{حالت ایده‌آل}

هرچند هدف اصلی ZSSR حالت ایده‌آل نیست، عملکرد آن بر مجموعه‌داده‌های
استاندارد ابرتفکیک ایده‌آل نیز ارزیابی شده است. در این
مجموعه‌داده‌ها، تصاویر LR با دستور \texttt{imresize} متلب (هستهٔ
دومکعبی با ضدتداخل) از نسخهٔ HR تولید شده‌اند.

\begin{table}[!ht]
	\centering
	\caption{مقایسهٔ نتایج ابرتفکیک در حالت ایده‌آل (کاهش‌نمونه‌برداری
		دومکعبی) بر حسب \lr{PSNR / SSIM}. برگرفته از \cite{shocher2018zssr}.}
	\label{tbl:zssr_ideal}
	\renewcommand{\arraystretch}{1.4}
	\begin{tabular}{|c|c|c|c|c|c|}
		\hline
		\textbf{مجموعه‌داده} & \textbf{مقیاس} &
		\textbf{SRCNN} & \textbf{VDSR} & \textbf{EDSR+} &
		\textbf{ZSSR} \\
		\hline
		\multirow{3}{*}{\lr{Set5}}
		& $\times 2$ & \lr{36.66/0.954} & \lr{37.53/0.959} &
		\lr{\textbf{38.20}/\textbf{0.961}} & \lr{37.37/0.957} \\
		& $\times 3$ & \lr{32.75/0.909} & \lr{33.66/0.921} &
		\lr{\textbf{34.76}/\textbf{0.929}} & \lr{33.42/0.919} \\
		& $\times 4$ & \lr{30.48/0.863} & \lr{31.35/0.884} &
		\lr{\textbf{32.62}/\textbf{0.898}} & \lr{31.13/0.880} \\
		\hline
		\multirow{3}{*}{\lr{Set14}}
		& $\times 2$ & \lr{32.42/0.906} & \lr{33.03/0.912} &
		\lr{\textbf{34.02}/\textbf{0.920}} & \lr{33.00/0.911} \\
		& $\times 3$ & \lr{29.28/0.821} & \lr{29.77/0.831} &
		\lr{\textbf{30.66}/\textbf{0.848}} & \lr{29.80/0.830} \\
		& $\times 4$ & \lr{27.49/0.750} & \lr{28.01/0.767} &
		\lr{\textbf{28.94}/\textbf{0.790}} & \lr{28.01/0.765} \\
		\hline
		\multirow{3}{*}{\lr{BSD100}}
		& $\times 2$ & \lr{31.36/0.888} & \lr{31.90/0.896} &
		\lr{\textbf{32.37}/\textbf{0.902}} & \lr{31.65/0.892} \\
		& $\times 3$ & \lr{28.41/0.786} & \lr{28.82/0.798} &
		\lr{\textbf{29.32}/\textbf{0.810}} & \lr{28.67/0.795} \\
		& $\times 4$ & \lr{26.90/0.710} & \lr{27.29/0.725} &
		\lr{\textbf{27.79}/\textbf{0.744}} & \lr{27.12/0.721} \\
		\hline
	\end{tabular}
\end{table}

نتایج جدول \ref{tbl:zssr_ideal} نشان می‌دهد که ZSSR خاصِ تصویر، در
مقابل روش‌های نظارت‌شده‌ای که به‌صورت گسترده برای همین شرایط آموزش
دیده‌اند، نتایج رقابتی کسب می‌کند. به‌طور خاص، ZSSR به‌مراتب بهتر از
SRCNN عمل می‌کند و در برخی موارد با
VDSR (که تا یک سال قبل پیشرفته‌ترین روش بود)
قابل مقایسه یا حتی بهتر است. در رژیم ابرتفکیک بدون نظارت نیز
ZSSR با اختلاف زیاد از روش پیشرو \lr{SelfExSR}
پیشی می‌گیرد.

نکتهٔ قابل توجه آن است که در تصاویری با ساختارهای تکرارشوندهٔ
درونی قوی، ZSSR حتی از VDSR و گاهی از EDSR+ نیز فراتر می‌رود،
حتی اگر تصاویر LR با تنظیمات ایده‌آل نظارت‌شده تولید شده باشند.
تحلیل بیشتر نشان می‌دهد که برخی پیکسل‌ها (به‌ویژه در نواحی با
تکرارپذیری بالای اطلاعات، مانند پنجره‌های بسیار کوچک در بالای
ساختمان) از یادگیری داخلی (ZSSR) بیشتر سود می‌برند، در حالی که
پیکسل‌های دیگر (مثلاً لبه‌هایی که در تصاویر طبیعی فراوان هستند)
از یادگیری خارجی (EDSR+) بهره‌مند‌ترند. این مشاهده حاکی از
پتانسیل ترکیب یادگیری داخلی و خارجی در یک چارچوب محاسباتی واحد
برای بهبود بیشتر عملکرد است.

\subsubsection{حالت غیرایده‌آل}

تصاویر LR واقعی معمولاً به‌صورت ایده‌آل تولید نمی‌شوند. ZSSR در دو
دستهٔ اصلی از شرایط غیرایده‌آل ارزیابی شده است:
(الف) هسته‌های کاهش‌نمونه‌برداری غیرایده‌آل و
(ب) تصاویر LR با کیفیت پایین.

\paragraph{هسته‌های کاهش‌نمونه‌برداری غیرایده‌آل.}
برای ایجاد یک مجموعه‌دادهٔ کنترل‌شده، از مجموعهٔ BSD100
استفاده شد. تصاویر HR با هسته‌های گاوسی تصادفی (با اندازهٔ معقول)
کاهش‌نمونه‌برداری شدند. ماتریس کوواریانس $\boldsymbol{\Sigma}$ هستهٔ
کاهش‌نمونه‌برداری هر تصویر به‌صورت تصادفی و مطابق رابطهٔ زیر انتخاب شد:
\begin{align}
	\lambda_1, \lambda_2 &\sim \mathcal{U}\!\left[0, \tfrac{s}{2}\right],
	\qquad \theta \sim \mathcal{U}[0, \pi],
	\nonumber\\[4pt]
	\boldsymbol{\Lambda} &= \mathrm{diag}(\lambda_1, \lambda_2),
	\qquad
	\mathbf{U} =
	\begin{bmatrix}
		\cos\theta & -\sin\theta \\
		\sin\theta & \cos\theta
	\end{bmatrix},
	\nonumber\\[4pt]
	\boldsymbol{\Sigma} &= \mathbf{U}\,\boldsymbol{\Lambda}\,\mathbf{U}^{\!\top},
	\label{eq:zssr_kernel}
\end{align}
که $s$ ضریب کاهش‌نمونه‌برداری HR-LR است. بنابراین هر تصویر LR با یک
هستهٔ تصادفی متفاوت کاهش‌نمونه‌برداری شده است.

\begin{table}[!ht]
	\centering
	\caption{ابرتفکیک در حضور هسته‌های کاهش‌نمونه‌برداری ناشناخته بر
		مجموعهٔ BSD100 ($\times 2$). برگرفته از \cite{shocher2018zssr}.}
	\label{tbl:zssr_kernel}
	\renewcommand{\arraystretch}{1.4}
	\begin{tabular}{|c|c|c|}
		\hline
		\textbf{روش} & \textbf{\lr{PSNR}} & \textbf{\lr{SSIM}} \\
		\hline
		\lr{VDSR} & \lr{27.72} & \lr{0.764} \\
		\lr{EDSR+} & \lr{27.78} & \lr{0.766} \\
		\lr{Blind-SR} & \lr{28.42} & \lr{0.783} \\
		\lr{ZSSR [هستهٔ تخمینی]} & \lr{28.81} & \lr{0.831} \\
		\lr{ZSSR [هستهٔ واقعی]} & \lr{\textbf{29.68}} & \lr{\textbf{0.841}} \\
		\hline
	\end{tabular}
\end{table}

نتایج جدول \ref{tbl:zssr_kernel} حاکی از آن است که ZSSR با اختلاف
بسیار زیاد از روش‌های پیشرفته پیشی می‌گیرد:
$+1$~dB برای هسته‌های ناشناخته (تخمین‌زده‌شده با \cite{michaeli2013blind})
و $+2$~dB هنگام ارائهٔ هسته‌های واقعی. از دیدگاه بصری نیز تصاویر SR
تولیدشده توسط روش‌های نظارت‌شده بسیار تار هستند، در حالی که ZSSR
جزئیات واضح‌تری بازسازی می‌کند.

نکتهٔ جالب آن است که روش بدون نظارت \lr{Blind-SR} \cite{michaeli2013blind}
که از یادگیری عمیق استفاده نمی‌کند نیز از روش‌های نظارت‌شدهٔ
پیشرفته بهتر عمل می‌کند. این یافته از تحلیل‌ها و مشاهدات قبلی پشتیبانی می‌کند مبنی بر اینکه
(الف) مدل دقیق کاهش‌نمونه‌برداری مهم‌تر از پیشین‌های پیچیدهٔ تصویر
است و (ب) استفاده از هستهٔ نادرست منجر به نتایج بیش‌هموار‌شده می‌گردد.

حالت خاص هستهٔ $\delta$ (منجر به تداخل‌گذاری\LTRfootnote{Aliasing})
نیز توسط روش‌های نظارت‌شده به‌خوبی مدیریت نمی‌شود، در حالی که ZSSR
قادر به مدیریت آن است.
\begin{figure}
	\centering
	\includegraphics[width=1\linewidth]{figures/ZSSR_nonideal.png}
	\caption{تصویری که با کرنل غیر ایده آل تخریب شده است. از چپ راست تصاویر نشان‌دهنده حقیقت مرجع، نتیجه EDSR+ و نتیجه ZSSR می‌باشند.}
	\label{fig:ZSSR_nonideal}
\end{figure}

\paragraph{تصاویر LR با کیفیت پایین.}
در این آزمایش، برای هر تصویر از مجموعهٔ BSD100
یک نوع تخریب تصادفی از میان سه نوع انتخاب شد:
(i) نویز گاوسی با $\sigma = 0.05$،
(ii) نویز لکه‌ای\LTRfootnote{Speckle noise} با $\sigma = 0.05$،
(iii) فشرده‌سازی \lr{JPEG} با کیفیت ۴۵.

\begin{table}[!ht]
	\centering
	\caption{ابرتفکیک در حضور تخریب ناشناختهٔ تصویر بر مجموعهٔ BSD100
		($\times 2$). برگرفته از \cite{shocher2018zssr}.}
	\label{tbl:zssr_degradation}
	\renewcommand{\arraystretch}{1.4}
	\begin{tabular}{|c|c|c|}
		\hline
		\textbf{روش} & \textbf{\lr{PSNR}} & \textbf{\lr{SSIM}} \\
		\hline
		درون‌یابی دومکعبی & \lr{27.92} & \lr{0.750} \\
		\lr{EDSR+} \cite{lim2017enhanceddeepresidualnetworks} & \lr{27.56} & \lr{0.714} \\
		\lr{ZSSR} & \lr{\textbf{28.61}} & \lr{\textbf{0.781}} \\
		\hline
	\end{tabular}
\end{table}

نتایج جدول \ref{tbl:zssr_degradation} نشان می‌دهد که ZSSR در مقابل
تخریب‌های ناشناخته مقاوم است، در حالی که روش‌های نظارت‌شدهٔ پیشرفته
به حدی آسیب می‌بینند که حتی درون‌یابی دومکعبی ساده از آن‌ها بهتر عمل
می‌کند. این یافته به‌وضوح شکنندگی روش‌های نظارت‌شده در شرایط واقعی
را نمایان می‌سازد.

\subsection{تحلیل مزایا و محدودیت‌ها}
\label{subsec:zssr_analysis}

مزایای کلیدی روش ZSSR را می‌توان به‌صورت زیر خلاصه کرد:

\begin{enumerate}
	\item \textbf{نخستین روش ابرتفکیک مبتنی بر CNN بدون نظارت:}
	ZSSR نشان داد که امکان بهره‌گیری از قدرت یادگیری عمیق بدون نیاز
	به هیچ دادهٔ آموزشی خارجی وجود دارد.
	\item \textbf{مدیریت شرایط غیرایده‌آل:} این روش قادر به مدیریت
	طیف وسیعی از شرایط تصویربرداری غیرایده‌آل از جمله هسته‌های
	کاهش‌نمونه‌برداری ناشناخته، نویز حسگر، فشرده‌سازی و تداخل‌گذاری
	است.
	\item \textbf{عدم نیاز به پیش‌آموزش:} شبکه در زمان آزمون آموزش
	می‌بیند و به منابع محاسباتی متوسط نیاز دارد.
	\item \textbf{انعطاف‌پذیری مقیاس:} قابل اعمال برای هر اندازه و
	نظریتاً هر نسبت ابعاد.
	\item \textbf{سازگاری در زمان آزمون:} قابلیت تنظیم پارامترها
	(هسته، نویز، بازافکنی) بر اساس ویژگی‌های تصویر آزمون.
	\item \textbf{عملکرد پیشرفته در شرایط غیرایده‌آل:} بهبود
	$1$--$2$~dB نسبت به روش‌های نظارت‌شدهٔ پیشرفته بر تصاویر غیرایده‌آل.
\end{enumerate}

در عین حال، ZSSR محدودیت‌هایی نیز دارد. در حالت ایده‌آل (هستهٔ
دومکعبی، بدون نویز)، عملکرد آن از روش‌های نظارت‌شدهٔ بسیار عمیق مانند
EDSR+ فاصله دارد، هرچند نتایج رقابتی هستند. علاوه بر آن، آموزش در
زمان آزمون هزینهٔ زمانی اضافی دارد، هرچند همان‌طور که بحث شد، این
هزینه برای تصاویر بزرگ قابل مقایسه یا حتی کمتر از زمان استنتاج
روش‌های بسیار عمیق است.

\subsection{نتیجه‌گیری و تأثیر بر پژوهش‌های بعدی}
\label{subsec:zssr_conclusion}

روش ZSSR \cite{shocher2018zssr} مفهوم ابرتفکیک «بدون نمونهٔ آموزشی»
را معرفی کرد که از قدرت یادگیری عمیق بدون اتکا به هیچ نمونهٔ خارجی
یا پیش‌آموزش بهره می‌گیرد. این امر از طریق یک شبکهٔ CNN کوچک خاص
تصویر محقق شد که در زمان آزمون بر نمونه‌های درونی استخراج‌شده صرفاً
از تصویر LR آزمون آموزش می‌بیند. این رویکرد ابرتفکیک تصاویر واقعی
را ممکن ساخت؛ تصاویری که فرایند تصویربرداری آن‌ها غیرایده‌آل، ناشناخته
و از تصویری به تصویر دیگر متفاوت است. در این شرایط غیرایده‌آل واقعی،
ZSSR هم از نظر کیفی و هم از نظر کمّی به‌طور قابل توجهی از روش‌های
پیشرفتهٔ ابرتفکیک پیشی گرفت.

مهم‌تر از نتایج کمّی، ZSSR پارادایم جدیدی را در ابرتفکیک تصویر
بنیان‌گذاری کرد: ترکیب قدرت تعمیم‌پذیری یادگیری عمیق با بهره‌گیری از
آمار درونی تصویر. این پارادایم مسیر را برای پژوهش‌های بعدی هموار
ساخت، از جمله روش‌هایی که یادگیری داخلی و خارجی را در یک چارچوب واحد
ترکیب می‌کنند و شبکه‌هایی که در زمان آزمون خود را با شرایط خاص
تصویر سازگار می‌سازند. این پژوهش همچنین اهمیت مدل دقیق فرایند
تخریب تصویر را بیش از پیشین‌های پیچیدهٔ تصویر برجسته ساخت؛ مشاهده‌ای
که تأثیر عمیقی بر جهت‌گیری تحقیقات ابرتفکیک در سال‌های بعد داشت.